%! Sinusoidal Function Transformations — Unit 6, Lesson 6.7 — Guided Notes
\documentclass[10pt]{article}
\usepackage{precalculus-article}
\usepackage{precalculus-boxes}
\newcommand{\TallMath}[1]{$\displaystyle #1\rule[-1.4em]{0pt}{3.2em}$}

\begin{document}

\pageheader{Unit 6, Lesson 6.7}{Guided Notes}
\namedateperiod

\vspace{0.10in}

%----------------------------------------------------------------------
% Objectives
%----------------------------------------------------------------------
\begin{objectivebox}
\textbf{I will be able to\ldots}
\begin{enumerate}[leftmargin=1.5em, itemsep=4pt]
  \item \textbf{LO 3.6.A} --- Identify the amplitude, vertical shift (midline),
    period, and phase shift of a sinusoidal function from its equation. \writeline
  \item \textbf{LO 3.6.A} --- Explain how each parameter in $f(\theta)=a\sin(b(\theta+c))+d$
    corresponds to an additive or multiplicative transformation of $y=\sin\theta$. \writeline
\end{enumerate}
\end{objectivebox}

\vspace{0.08in}

%----------------------------------------------------------------------
% Vocabulary
%----------------------------------------------------------------------
\begin{vocabbox}
\begin{description}[leftmargin=0pt, itemsep=6pt]
  \item[\termblanklong{sinusoidal function}]
    A function equivalent to $a\sin(b(\theta+c))+d$ (or the cosine form) with $a \neq 0$.
    A \underline{\hspace{3cm}} of the parent sine or cosine function.

  \item[\termblanklong{amplitude}]
    Equals $|a|$ --- half the vertical distance from the
    \underline{\hspace{2cm}} to the \underline{\hspace{2cm}}.
    Controls the vertical \underline{\hspace{2.5cm}}.

  \item[\termblanklong{midline}]
    The horizontal line $y = $ \underline{\hspace{1cm}}. The vertical
    \underline{\hspace{2cm}} of the function. Average of max and min outputs.

  \item[\termblanklong{period}]
    Length of one complete cycle; $T = \dfrac{2\pi}{\underline{\hspace{0.8cm}}}$.
    Controlled by the horizontal dilation factor $b$.

  \item[\termblanklong{phase shift}]
    Horizontal translation equal to \underline{\hspace{1.5cm}} when the function
    is written in factored form $b(\theta + c)$.
    Positive $c$ shifts \underline{\hspace{1.2cm}}; negative $c$ shifts
    \underline{\hspace{1.2cm}}.
\end{description}
\end{vocabbox}

\vspace{0.08in}

%----------------------------------------------------------------------
% Hook
%----------------------------------------------------------------------
\begin{hookbox}
\textbf{Entry Question:} The two sinusoidal functions below are shown on
the same axes. Use the graphs to answer the questions.

\vspace{4pt}
\begin{center}
\begin{tikzpicture}[x=1.0cm, y=0.9cm]
  \draw[linegray, very thin] (-0.2,-1.6) grid[xstep=0.5,ystep=0.5] (6.7,3.6);
  \draw[->] (-0.2,0) -- (6.8,0) node[right]{\small $\theta$};
  \draw[->] (0,-1.6) -- (0,3.7) node[above]{\small $y$};
  \foreach \x/\lbl in {1.571/$\frac{\pi}{2}$, 3.142/$\pi$, 4.712/$\frac{3\pi}{2}$, 6.283/$2\pi$} {
    \draw (\x,0.07)--(\x,-0.07);
    \node[below,font=\scriptsize] at (\x,-0.07) {\lbl};
  }
  \foreach \y/\lbl in {-1/{$-1$},1/{$1$},2/{$2$},3/{$3$}} {
    \draw (0.07,\y)--(-0.07,\y);
    \node[left,font=\scriptsize] at (-0.07,\y) {\lbl};
  }
  % Graph A: y = sin(theta)
  \draw[navy, thick, domain=0:6.3, samples=100]
    plot (\x, {sin(deg(\x))}) node[right,font=\small]{\textbf{A}};
  % Graph B: y = 3sin(theta)+1
  \draw[redacc, thick, domain=0:6.3, samples=100]
    plot (\x, {3*sin(deg(\x))+1}) node[right,font=\small]{\textbf{B}};
\end{tikzpicture}
\end{center}

\vspace{4pt}
\textbf{Q1.} Graph A has amplitude \underline{\hspace{1.5cm}} and midline $y = $ \underline{\hspace{1.5cm}}.\\[4pt]
\textbf{Q2.} Graph B has amplitude \underline{\hspace{1.5cm}} and midline $y = $ \underline{\hspace{1.5cm}}.\\[4pt]
\textbf{Q3.} What equation could describe Graph B?
$y = $ \underline{\hspace{3.5cm}}

\end{hookbox}

\vspace{0.08in}

%----------------------------------------------------------------------
% LO 3.6.A (Part 1) — Vertical Transformations
%----------------------------------------------------------------------
\begin{notesbox}{LO 3.6.A --- Part 1: Vertical Transformations ($a$ and $d$)}

\textbf{Vertical Dilation:} $g(\theta) = a\sin(\theta)$

Starting from $y = \sin\theta$, multiplying all outputs by $a$ changes the
\underline{\hspace{2cm}}. The new amplitude is \underline{\hspace{1.5cm}}.
If $a < 0$, the graph is also \underline{\hspace{3cm}} over the midline.

\vspace{6pt}
\renewcommand{\arraystretch}{1.5}
\begin{tabular}{|l|c|c|c|c|}
\hline
\textbf{Function} & \textbf{Amplitude} & \textbf{Max} & \textbf{Min} & \textbf{Midline} \\
\hline
$y = \sin\theta$ & & & & \\
\hline
$y = 4\sin\theta$ & & & & \\
\hline
$y = -2\sin\theta$ & & & & \\
\hline
$y = \tfrac{1}{2}\sin\theta$ & & & & \\
\hline
\end{tabular}

\vspace{8pt}
\textbf{Vertical Translation:} $g(\theta) = \sin(\theta) + d$

Adding $d$ to all outputs shifts the \underline{\hspace{2cm}} to $y = d$.
The amplitude is \underline{\hspace{3cm}} (unchanged).

\vspace{6pt}
\begin{tabular}{|l|c|c|c|c|}
\hline
\textbf{Function} & \textbf{Amplitude} & \textbf{Max} & \textbf{Min} & \textbf{Midline} \\
\hline
$y = \sin\theta + 3$ & & & & \\
\hline
$y = 2\sin\theta - 1$ & & & & \\
\hline
\end{tabular}

\end{notesbox}

\vspace{0.08in}

%----------------------------------------------------------------------
% LO 3.6.A (Part 2) — Horizontal Transformations
%----------------------------------------------------------------------
\begin{notesbox}{LO 3.6.A --- Part 2: Horizontal Transformations ($b$ and $c$)}

\textbf{Horizontal Dilation:} $g(\theta) = \sin(b\theta)$

Replacing $\theta$ with $b\theta$ compresses or stretches the graph horizontally.
New period $= \dfrac{2\pi}{|b|}$.

\vspace{4pt}
If $|b| > 1$: the graph is \underline{\hspace{2.5cm}} (period is shorter). \\
If $|b| < 1$: the graph is \underline{\hspace{2.5cm}} (period is longer).

\vspace{6pt}
\begin{tabular}{|l|c|c|}
\hline
\textbf{Function} & \textbf{$b$ value} & \textbf{Period} \\
\hline
$y = \sin(3\theta)$ & & \\
\hline
$y = \sin\!\left(\tfrac{1}{2}\theta\right)$ & & \\
\hline
$y = \sin(4\theta)$ & & \\
\hline
\end{tabular}

\vspace{8pt}
\textbf{Phase Shift (Horizontal Translation):} $g(\theta) = \sin(\theta + c)$

Replacing $\theta$ with $\theta + c$ shifts the graph \underline{\hspace{1cm}}
by $|c|$ units.
Phase shift $= $ \underline{\hspace{1.5cm}} (use the factored form to read it correctly).

\vspace{4pt}
\textbf{Important --- Factoring first!} Before reading the phase shift, factor:
$$\sin(2\theta + \pi) = \sin\!\bigl(2(\theta + \underline{\hspace{1cm}})\bigr)$$
Phase shift $= -\underline{\hspace{1cm}} = \underline{\hspace{1.5cm}}$

\end{notesbox}

\vspace{0.08in}

%----------------------------------------------------------------------
% LO 3.6.A (Part 3) — General Form and Summary
%----------------------------------------------------------------------
\begin{notesbox}{LO 3.6.A --- Part 3: The General Form}

\textbf{General Sinusoidal Function:}
$$f(\theta) = a\sin\bigl(b(\theta + c)\bigr) + d$$

\renewcommand{\arraystretch}{1.6}
\begin{tabularx}{\linewidth}{@{} p{0.14\linewidth} p{0.26\linewidth} X @{}}
\hline
\textbf{Parameter} & \textbf{Feature} & \textbf{Formula} \\
\hline
$a$ & Amplitude & $|a|$ \quad (sign determines reflection) \\
$b$ & Period & $T = \dfrac{2\pi}{|b|}$ \\
$c$ & Phase shift & $-c$ \quad (left if $c > 0$, right if $c < 0$) \\
$d$ & Midline & $y = d$ \\
\hline
\end{tabularx}

\vspace{8pt}
\textbf{Worked Example:} Identify all features of
$f(\theta) = -2\sin\!\bigl(\tfrac{1}{2}(\theta - \pi)\bigr) + 3$.

\vspace{4pt}
\begin{tabular}{ll}
Amplitude: & \underline{\hspace{3cm}} \\[6pt]
Period: & \underline{\hspace{3cm}} \\[6pt]
Phase shift: & \underline{\hspace{3cm}} \\[6pt]
Midline: & \underline{\hspace{3cm}} \\[6pt]
Reflected? & \underline{\hspace{3cm}} \\
\end{tabular}

\vspace{8pt}
\textbf{Sine--Cosine Equivalence (EK 3.6.A.1):}
$$\cos\theta = \sin\!\left(\theta + \dfrac{\pi}{2}\right)$$
Any cosine function can be rewritten as a sine function with a phase shift.

\vspace{4pt}
\textbf{Example:} Rewrite $2\cos\theta$ as a sine function.
$$2\cos\theta = 2\sin\!\left(\theta + \underline{\hspace{1.5cm}}\right)$$

\end{notesbox}

\vspace{0.08in}

%----------------------------------------------------------------------
% Practice
%----------------------------------------------------------------------
\begin{practicebox}

\textbf{Guided Practice 1.}\quad For each function, identify the amplitude,
period, phase shift, and midline.

\vspace{4pt}
\renewcommand{\arraystretch}{1.8}
\begin{tabular}{|l|c|c|c|c|}
\hline
\textbf{Function} & \textbf{Amplitude} & \textbf{Period} & \textbf{Phase Shift} & \textbf{Midline} \\
\hline
$f(\theta) = 5\sin(3\theta) - 2$ & & & & \\
\hline
$g(\theta) = -\cos\!\left(\tfrac{\pi}{2}\theta\right) + 4$ & & & & \\
\hline
$h(\theta) = 3\sin\!\left(2\left(\theta + \dfrac{\pi}{6}\right)\right)$ & & & & \\
\hline
\end{tabular}

\vspace{0.10in}
\textbf{Guided Practice 2.}\quad The function $k(\theta) = 4\sin(3\theta - \pi) + 1$.

\vspace{4pt}
(a)\quad Factor the argument: $3\theta - \pi = 3\!\left(\theta - \underline{\hspace{1.5cm}}\right)$

\vspace{4pt}
(b)\quad Phase shift $= $ \underline{\hspace{2.5cm}} (left or right by how much?)

\vspace{4pt}
(c)\quad List all four features: amplitude \underline{\hspace{1.5cm}},
period \underline{\hspace{1.5cm}}, phase shift \underline{\hspace{1.5cm}},
midline \underline{\hspace{1.5cm}}.

\end{practicebox}

\end{document}
